\documentclass[11pt,a4paper]{article}
\usepackage[margin=19mm,top=17mm,bottom=17mm]{geometry}
\usepackage{amsmath,amssymb,amsfonts}
\usepackage{mathptmx}
\usepackage{enumitem}
\usepackage{fancyhdr}
\usepackage{array,tabularx,booktabs}
\usepackage[hidelinks]{hyperref}
\setlength{\parindent}{0pt}\setlength{\parskip}{4pt}
\setlist[enumerate,1]{leftmargin=1.1em,itemsep=3pt,parsep=2pt}
\setlist[itemize]{leftmargin=1.4em,itemsep=1pt}
\pagestyle{fancy}\fancyhf{}
\fancyhead[C]{\footnotesize \textbf{Kinematics test - mixed topic}}
\fancyfoot[C]{\footnotesize Page \thepage}
\renewcommand{\headrulewidth}{0.4pt}
\begin{document}
{\centering \large \textbf{Bodha Academy}\par}
{\centering \small \textit{Sector V, Kolkata}\par}
\vspace{2mm}
\noindent \textbf{Marks : 138} \hfill \textbf{Time : 2 : 00 hours}
\noindent\rule{\linewidth}{1.1pt}\vspace{-1mm}\noindent\rule{\linewidth}{0.5pt}
{\centering \Large \textbf{KINEMATICS TEST - MIXED TOPIC}\par}
{\centering JEE $|$ Physics\par}
\noindent\rule{\linewidth}{1.1pt}\vspace{-1mm}\noindent\rule{\linewidth}{0.5pt}
\vspace{1mm}
\noindent \textbf{Duration:} 2 : 00 hours \quad \textbf{Total Marks:} 138 \quad \textbf{Questions:} 40\par
\vspace{2mm}

\begin{tabularx}{\linewidth}{|X|X|}\hline
Candidate Name: \rule{6cm}{0.4pt} & Roll Number: \rule{3.5cm}{0.4pt} \\\hline
\end{tabularx}

\textbf{Instructions}
\begin{itemize}[topsep=2pt]
\item Read every question carefully before selecting an answer.
\item For single-correct questions, mark one option only.
\item Use the answer sheet or space provided by the invigilator.
\end{itemize}
\section*{Kinematics in One Dimension › Equations of uniformly accelerated motion}
\begin{enumerate}[resume]
\item Two particles move along the same straight track. The positive $x$-direction is to the right. At $t=0$, particle $A$ is at $x=0$ with velocity $2\,\text{m s}^{-1}$ and constant acceleration $1\,\text{m s}^{-2}$, while particle $B$ is at $x=27\,\text{m}$ with velocity $-4\,\text{m s}^{-1}$ and constant acceleration $-1\,\text{m s}^{-2}$. They start simultaneously. The first positive time at which they meet is \hfill [1 marks]
{\renewcommand{\arraystretch}{1.8}
\begin{tabularx}{\linewidth}{X X X X}
(A) \quad $2\,\text{s}$ & (B) \quad $3\,\text{s}$ & (C) \quad $4\,\text{s}$ & (D) \quad $6\,\text{s}$ \\[10pt]
\end{tabularx}}
\item B vehicle moves along a straight, level track with constant acceleration and never reverses its direction. A motion tracker records the average velocity of the vehicle during three consecutive time intervals, each of duration $T$. The average velocities during the first and second intervals are $8\,\mathrm{m\,s^{-1}}$ and $14\,\mathrm{m\,s^{-1}}$, respectively. If $s_1$ and $s_3$ denote the displacements during the first and third intervals, then the ratio $s_3/s_1$ is \hfill [4 marks]
{\renewcommand{\arraystretch}{1.8}
\begin{tabularx}{\linewidth}{X X X X}
(A) \quad $\dfrac{3}{2}$ & (B) \quad $2$ & (C) \quad $\dfrac{5}{2}$ & (D) \quad $3$ \\[10pt]
\end{tabularx}}
\item Two carts A and B move on a straight, frictionless track. At $t=0$, cart A is at $x=0$ and cart B is at $x=18\,\text{m}$. Both carts are released simultaneously from rest. Cart A has a constant acceleration of $2\,\text{m s}^{-2}$ toward B, while cart B has a constant acceleration of $1\,\text{m s}^{-2}$ toward A. Their positions are measured from A's initial position. At what position do the carts meet? \hfill [4 marks]
{\renewcommand{\arraystretch}{1.8}
\begin{tabularx}{\linewidth}{X X X X}
(A) \quad $6\,\text{m}$ & (B) \quad $12\,\text{m}$ & (C) \quad $18\,\text{m}$ & (D) \quad $\sqrt{12}\,\text{m}$ \\[10pt]
\end{tabularx}}
\item Particle A starts from a marker at time $t=0$ with initial velocity $6\,\text{m s}^{-1}$ and constant acceleration $2\,\text{m s}^{-2}$. Particle B starts from the same marker $2\,\text{s}$ later with initial velocity $10\,\text{m s}^{-1}$ and constant acceleration $4\,\text{m s}^{-2}$. Both particles move along the same straight track in the same direction. The time elapsed after particle B starts when they meet for the first time is \hfill [4 marks]
{\renewcommand{\arraystretch}{1.8}
\begin{tabularx}{\linewidth}{X X X X}
(A) \quad $3\,\text{s}$ & (B) \quad $4\,\text{s}$ & (C) \quad $5\,\text{s}$ & (D) \quad $6\,\text{s}$ \\[10pt]
\end{tabularx}}
\item Two cyclists move along the same straight road in the same direction. Cyclist A passes a fixed marker at time $t=0$ with speed $6\,\mathrm{m\,s^{-1}}$ and thereafter moves with constant acceleration $2\,\mathrm{m\,s^{-2}}$. Cyclist B passes the same marker $3\,\mathrm{s}$ later, starting from rest, and moves with constant acceleration $6\,\mathrm{m\,s^{-2}}$. The time elapsed after cyclist B starts when B catches A is \hfill [4 marks]
{\renewcommand{\arraystretch}{1.8}
\begin{tabularx}{\linewidth}{X X}
(A) \quad $\left(3+\frac{3\sqrt{10}}{2}\right)\,\mathrm{s}$ & (B) \quad $\left(3-\frac{3\sqrt{10}}{2}\right)\,\mathrm{s}$ \\[10pt]
(C) \quad $\left(6+\frac{3\sqrt{10}}{2}\right)\,\mathrm{s}$ & (D) \quad $\left(6-\frac{3\sqrt{10}}{2}\right)\,\mathrm{s}$ \\[10pt]
\end{tabularx}}
\item Two particles A and B move simultaneously along the same straight track in the same direction. At t = 0, particle A is 24 m behind particle B. Their initial velocities are 10 m/s and 15 m/s, respectively, while their constant accelerations are 5 m/s^2 and 3 m/s^2, respectively. The first time at which A catches B is: \hfill [4 marks]
{\renewcommand{\arraystretch}{1.8}
\begin{tabularx}{\linewidth}{X X X X}
(A) \quad 4 s & (B) \quad 6 s & (C) \quad 8 s & (D) \quad 10 s \\[10pt]
\end{tabularx}}
\item A particle moves along the x-axis with constant acceleration. It passes through the origin at t = 0, is at x = 12\,\text{m} at t = 2\,\text{s}, and passes through the origin again at t = 6\,\text{s}. What is the speed of the particle when it passes through x = 12\,\text{m} for the second time? \hfill [4 marks]
{\renewcommand{\arraystretch}{1.8}
\begin{tabularx}{\linewidth}{X X X X}
(A) \quad 3\,\text{m/s} & (B) \quad 6\,\text{m/s} & (C) \quad -3\,\text{m/s} & (D) \quad 9\,\text{m/s} \\[10pt]
\end{tabularx}}
\item Two particles A and B move along the same straight track. Take the positive direction from A toward B. At time $t=0$, particle A passes the marker $x=0$ with velocity $4\,\mathrm{m\,s^{-1}}$ and constant acceleration $2\,\mathrm{m\,s^{-2}}$. At the same instant, particle B is at the marker $x=56\,\mathrm{m}$, is released from rest, and moves toward A with constant acceleration $3\,\mathrm{m\,s^{-2}}$. The time after $t=0$ at which they meet is: \hfill [3 marks]
{\renewcommand{\arraystretch}{1.8}
\begin{tabularx}{\linewidth}{X X X X}
(A) \quad $3.8\,\mathrm{s}$ & (B) \quad $3.9\,\mathrm{s}$ & (C) \quad $4.0\,\mathrm{s}$ & (D) \quad $4.1\,\mathrm{s}$ \\[10pt]
\end{tabularx}}
\item Two automated carts move along the same straight track in the positive $x$-direction. Cart A passes the reference mark $x=0$ at $t=0$ with initial velocity $10\,\mathrm{m\,s^{-1}}$ and constant acceleration $4\,\mathrm{m\,s^{-2}}$. Cart B is initially at rest at the point $x=71\,\mathrm{m}$ and is released at $t=3\,\mathrm{s}$, after which it moves with constant acceleration $2\,\mathrm{m\,s^{-2}}$. The time, measured from $t=0$, at which Cart A catches Cart B is \hfill [4 marks]
{\renewcommand{\arraystretch}{1.8}
\begin{tabularx}{\linewidth}{X X X X}
(A) \quad $2\,\mathrm{s}$ & (B) \quad $3\,\mathrm{s}$ & (C) \quad $4\,\mathrm{s}$ & (D) \quad $5\,\mathrm{s}$ \\[10pt]
\end{tabularx}}
\item A train enters a $120\,\mathrm{m}$ long platform with a speed of $18\,\mathrm{m\,s^{-1}}$. It is uniformly decelerated from the instant it enters the platform and comes to rest exactly at the far end of the platform. A warning signal is received $4\,\mathrm{s}$ after the train enters the platform. The time after receiving the warning signal at which the train comes to rest is \hfill [4 marks]
{\renewcommand{\arraystretch}{1.8}
\begin{tabularx}{\linewidth}{X X}
(A) \quad $\dfrac{28}{3}\,\mathrm{s}$ & (B) \quad $\dfrac{40}{3}\,\mathrm{s}$ \\[10pt]
(C) \quad $\dfrac{52}{3}\,\mathrm{s}$ & (D) \quad $4\,\mathrm{s}$ \\[10pt]
\end{tabularx}}
\end{enumerate}
\section*{Kinematics in One Dimension › Motion graphs}
\begin{enumerate}[resume]
\item A particle moves along the $x$-axis. Its velocity varies with time as follows: from $t=0$ to $t=2\,\text{s}$, the velocity decreases uniformly from $4\,\text{m s}^{-1}$ to $0$; from $t=2$ to $t=4\,\text{s}$, it decreases uniformly from $0$ to $-2\,\text{m s}^{-1}$; and from $t=4$ to $t=6\,\text{s}$, it remains constant at $-2\,\text{m s}^{-1}$. The displacement of the particle during the first $6\,\text{s}$ is \hfill [4 marks]
{\renewcommand{\arraystretch}{1.8}
\begin{tabularx}{\linewidth}{X X X X}
(A) \quad $-2\,\text{m}$ & (B) \quad $2\,\text{m}$ & (C) \quad $6\,\text{m}$ & (D) \quad $10\,\text{m}$ \\[10pt]
\end{tabularx}}
\item A particle moves along the $x$-axis. Its velocity varies linearly with time as follows: the velocity increases uniformly from $0$ to $4\,\text{m s}^{-1}$ during the first $2\,\text{s}$, decreases uniformly to $-2\,\text{m s}^{-1}$ during the next $2\,\text{s}$, and then remains constant at $-2\,\text{m s}^{-1}$ for the next $2\,\text{s}$. The displacement of the particle during the complete $6\,\text{s}$ interval is \hfill [4 marks]
{\renewcommand{\arraystretch}{1.8}
\begin{tabularx}{\linewidth}{X X X X}
(A) \quad $-2\,\text{m}$ & (B) \quad $0\,\text{m}$ & (C) \quad $2\,\text{m}$ & (D) \quad $6\,\text{m}$ \\[10pt]
\end{tabularx}}
\item A particle moves along a straight line. Its velocity–time graph consists of the following three straight-line segments: from $t=0$ to $t=2\,\text{s}$, velocity increases uniformly from $0$ to $4\,\text{m s}^{-1}$; from $t=2$ to $t=4\,\text{s}$, velocity remains constant at $4\,\text{m s}^{-1}$; and from $t=4$ to $t=6\,\text{s}$, velocity decreases uniformly from $4$ to $-4\,\text{m s}^{-1}$. The displacement and total distance travelled during the first $6\,\text{s}$ are, respectively: \hfill [4 marks]
{\renewcommand{\arraystretch}{1.8}
\begin{tabularx}{\linewidth}{X X}
(A) \quad $12\,\text{m},\ 12\,\text{m}$ & (B) \quad $12\,\text{m},\ 16\,\text{m}$ \\[10pt]
(C) \quad $16\,\text{m},\ 12\,\text{m}$ & (D) \quad $16\,\text{m},\ 16\,\text{m}$ \\[10pt]
\end{tabularx}}
\item A particle moves along the $x$-axis. Its velocity varies with time as follows: from $t=0$ to $t=4\,\text{s}$, the velocity decreases uniformly from $4\,\text{m s}^{-1}$ to $-2\,\text{m s}^{-1}$; from $t=4$ to $t=6\,\text{s}$, it remains constant at $-2\,\text{m s}^{-1}$; and from $t=6$ to $t=8\,\text{s}$, it increases uniformly to zero. The displacement and the total distance travelled by the particle during the first $8\,\text{s}$ are, respectively: \hfill [4 marks]
{\renewcommand{\arraystretch}{1.8}
\begin{tabularx}{\linewidth}{X X}
(A) \quad $-2\,\text{m},\ \dfrac{38}{3}\,\text{m}$ & (B) \quad $2\,\text{m},\ \dfrac{38}{3}\,\text{m}$ \\[10pt]
(C) \quad $-2\,\text{m},\ \dfrac{22}{3}\,\text{m}$ & (D) \quad $2\,\text{m},\ \dfrac{22}{3}\,\text{m}$ \\[10pt]
\end{tabularx}}
\item The velocity–time graph of a particle moving along a straight line is formed by joining the points $(0,0)$, $(2,4)$, $(4,4)$ and $(6,-2)$ by straight-line segments, where time is measured in seconds and velocity in $\text{m s}^{-1}$. The displacement and the total distance travelled by the particle during the first $6$ s are, respectively, \hfill [4 marks]
{\renewcommand{\arraystretch}{1.8}
\begin{tabularx}{\linewidth}{X X}
(A) \quad $14\ \text{m},\ \dfrac{46}{3}\ \text{m}$ & (B) \quad $\dfrac{46}{3}\ \text{m},\ 14\ \text{m}$ \\[10pt]
(C) \quad $14\ \text{m},\ \dfrac{40}{3}\ \text{m}$ & (D) \quad $\dfrac{40}{3}\ \text{m},\ 14\ \text{m}$ \\[10pt]
\end{tabularx}}
\item A particle moves along the $x$-axis. Its velocity varies with time as follows: it increases uniformly from $0$ to $4\,\text{m s}^{-1}$ during the first $2\,\text{s}$, remains constant at $4\,\text{m s}^{-1}$ from $t=2\,\text{s}$ to $t=4\,\text{s}$, and then decreases uniformly to $-2\,\text{m s}^{-1}$ at $t=6\,\text{s}$. It remains at $-2\,\text{m s}^{-1}$ from $t=6\,\text{s}$ to $t=8\,\text{s}$. The displacement of the particle during the first $8\,\text{s}$ is \hfill [1 marks]
{\renewcommand{\arraystretch}{1.8}
\begin{tabularx}{\linewidth}{X X X X}
(A) \quad $6\,\text{m}$ & (B) \quad $8\,\text{m}$ & (C) \quad $10\,\text{m}$ & (D) \quad $12\,\text{m}$ \\[10pt]
\end{tabularx}}
\item The velocity–time graph of a particle moving along a straight line is described as follows: its velocity increases uniformly from $0$ to $4\,\text{m s}^{-1}$ during the first $2\,\text{s}$, remains constant at $4\,\text{m s}^{-1}$ from $t=2\,\text{s}$ to $t=4\,\text{s}$, and then decreases uniformly to $-4\,\text{m s}^{-1}$ at $t=6\,\text{s}$. The displacement and the total distance travelled by the particle in the first $6\,\text{s}$ are, respectively: \hfill [4 marks]
{\renewcommand{\arraystretch}{1.8}
\begin{tabularx}{\linewidth}{X X}
(A) \quad $12\,\text{m},\ 16\,\text{m}$ & (B) \quad $16\,\text{m},\ 12\,\text{m}$ \\[10pt]
(C) \quad $8\,\text{m},\ 16\,\text{m}$ & (D) \quad $12\,\text{m},\ 12\,\text{m}$ \\[10pt]
\end{tabularx}}
\item The velocity–time graph of a particle moving along a straight line is described as follows: its velocity increases uniformly from $0$ to $4\,\text{m s}^{-1}$ during the first $2\,\text{s}$, remains constant at $4\,\text{m s}^{-1}$ from $t=2\,\text{s}$ to $t=4\,\text{s}$, and then decreases uniformly to $-2\,\text{m s}^{-1}$ at $t=6\,\text{s}$. The displacement and distance travelled by the particle in the first $6\,\text{s}$ are, respectively: \hfill [4 marks]
{\renewcommand{\arraystretch}{1.8}
\begin{tabularx}{\linewidth}{X X}
(A) \quad $14\,\text{m},\ \dfrac{46}{3}\,\text{m}$ & (B) \quad $14\,\text{m},\ \dfrac{38}{3}\,\text{m}$ \\[10pt]
(C) \quad $\dfrac{46}{3}\,\text{m},\ 14\,\text{m}$ & (D) \quad $12\,\text{m},\ 16\,\text{m}$ \\[10pt]
\end{tabularx}}
\item A particle moves along the x-axis. Its velocity–time graph is described as follows: from t=0 to t=2 s, its velocity increases uniformly from 0 to 6 m/s; from t=2 s to t=4 s, the velocity remains constant at 6 m/s; and from t=4 s to t=6 s, the velocity decreases uniformly from 6 m/s to -2 m/s. The displacement and total distance travelled by the particle during the first 6 s are, respectively, \hfill [1 marks]
{\renewcommand{\arraystretch}{1.8}
\begin{tabularx}{\linewidth}{X X X X}
(A) \quad 22 m and 23 m & (B) \quad 23 m and 22 m & (C) \quad 20 m and 24 m & (D) \quad 24 m and 24 m \\[10pt]
\end{tabularx}}
\item The velocity–time graph of a particle moving along a straight line is formed by joining, in order, the points $(0,0)$, $(2,4)$, $(5,4)$, $(7,-4)$ and $(9,0)$ by straight-line segments. Here, time is measured in seconds and velocity in $\text{m s}^{-1}$. The displacement and the total distance travelled by the particle during the first $9$ seconds are, respectively, \hfill [4 marks]
{\renewcommand{\arraystretch}{1.8}
\begin{tabularx}{\linewidth}{X X}
(A) \quad $12\,\text{m},\ 24\,\text{m}$ & (B) \quad $24\,\text{m},\ 12\,\text{m}$ \\[10pt]
(C) \quad $16\,\text{m},\ 24\,\text{m}$ & (D) \quad $12\,\text{m},\ 20\,\text{m}$ \\[10pt]
\end{tabularx}}
\end{enumerate}
\section*{Kinematics in One Dimension › Position, velocity and acceleration}
\begin{enumerate}[resume]
\item A particle moves along a straight track with velocity $v(t)=t^2-4t+3\ \text{m s}^{-1}$ for $0\le t\le 5\ \text{s}$. If its position at $t=0$ is $x(0)=0$, what is the total distance travelled by the particle during this interval? \hfill [4 marks]
{\renewcommand{\arraystretch}{1.8}
\begin{tabularx}{\linewidth}{X X}
(A) \quad $\frac{20}{3}\ \text{m}$ & (B) \quad $4\ \text{m}$ \\[10pt]
(C) \quad $\frac{28}{3}\ \text{m}$ & (D) \quad $12\ \text{m}$ \\[10pt]
\end{tabularx}}
\item A particle moves along a straight line with velocity $v(t)=(6-2t)\,\mathrm{m\,s^{-1}}$ for $0\le t\le 5\,\mathrm{s}$. It starts from position $x=0$ at $t=0$. The ratio of the total distance travelled by the particle to the magnitude of its net displacement during this interval is \hfill [1 marks]
{\renewcommand{\arraystretch}{1.8}
\begin{tabularx}{\linewidth}{X X X X}
(A) \quad $\dfrac{5}{13}$ & (B) \quad $\dfrac{13}{5}$ & (C) \quad $\dfrac{9}{4}$ & (D) \quad $\dfrac{13}{9}$ \\[10pt]
\end{tabularx}}
\item A particle moves along a straight track. Its velocity is constant at $+4\,\mathrm{m\,s^{-1}}$ from $t=0$ to $t=3\,\mathrm{s}$. From $t=3\,\mathrm{s}$ to $t=5\,\mathrm{s}$, its velocity changes uniformly from $+4\,\mathrm{m\,s^{-1}}$ to $-2\,\mathrm{m\,s^{-1}}$, and it remains constant thereafter. If $x(0)=0$, which one of the following statements is correct? \hfill [1 marks]
\begin{enumerate}[label=(\Alph*),leftmargin=1.6em,itemsep=1pt]\item At $t=5\,\mathrm{s}$, $x=10\,\mathrm{m}$, and the particle has never reversed its direction. \item At $t=5\,\mathrm{s}$, $x=14\,\mathrm{m}$, but the particle is still moving in the positive direction. \item At $t=5\,\mathrm{s}$, $x=14\,\mathrm{m}$, the velocity is $-2\,\mathrm{m\,s^{-1}}$, and the particle has reversed its direction exactly once. \item At $t=5\,\mathrm{s}$, $x=16\,\mathrm{m}$, and the particle has reversed its direction twice.\end{enumerate}
\item A particle moves along the $x$-axis with velocity $v(t)=t^2-4t+3\ \mathrm{m\,s^{-1}}$ for $0\le t\le 5\ \mathrm{s}$. If its position at $t=0$ is $x(0)=0$, then the total distance travelled by the particle during the first $5$ seconds is \hfill [4 marks]
{\renewcommand{\arraystretch}{1.8}
\begin{tabularx}{\linewidth}{X X}
(A) \quad $\dfrac{20}{3}\ \mathrm{m}$ & (B) \quad $\dfrac{28}{3}\ \mathrm{m}$ \\[10pt]
(C) \quad $\dfrac{32}{3}\ \mathrm{m}$ & (D) \quad $12\ \mathrm{m}$ \\[10pt]
\end{tabularx}}
\item A particle starts from rest at the origin and moves along the positive x-axis. It accelerates uniformly at $2\,\mathrm{m\,s^{-2}}$ for $5\,\mathrm{s}$. It then continues with the velocity attained for $4\,\mathrm{s}$. Finally, it decelerates uniformly at $5\,\mathrm{m\,s^{-2}}$ until it comes to rest, without reversing its direction. Which of the following correctly gives the total displacement from the origin and the duration of the final deceleration stage, respectively? \hfill [4 marks]
{\renewcommand{\arraystretch}{1.8}
\begin{tabularx}{\linewidth}{X X}
(A) \quad $65\,\mathrm{m},\ 2\,\mathrm{s}$ & (B) \quad $75\,\mathrm{m},\ 2\,\mathrm{s}$ \\[10pt]
(C) \quad $75\,\mathrm{m},\ 4\,\mathrm{s}$ & (D) \quad $85\,\mathrm{m},\ 2\,\mathrm{s}$ \\[10pt]
\end{tabularx}}
\item A particle moves along the positive $x$-axis. Its velocity is zero at $t=0$ and $t=4\,\text{s}$, and its velocity remains positive for $0<t<4\,\text{s}$. The acceleration is constant during each of the intervals $0\le t\le2\,\text{s}$ and $2\le t\le4\,\text{s}$, but changes suddenly at $t=2\,\text{s}$. The displacement during the first $2\,\text{s}$ is $12\,\text{m}$, while the total displacement during the $4\,\text{s}$ is $24\,\text{m}$. The acceleration during the interval $2\le t\le4\,\text{s}$ is: \hfill [4 marks]
{\renewcommand{\arraystretch}{1.8}
\begin{tabularx}{\linewidth}{X X X X}
(A) \quad $-3\,\text{m s}^{-2}$ & (B) \quad $-6\,\text{m s}^{-2}$ & (C) \quad $+6\,\text{m s}^{-2}$ & (D) \quad $-12\,\text{m s}^{-2}$ \\[10pt]
\end{tabularx}}
\item A particle starts from the origin and moves along the positive x-axis with initial velocity \(6\,\text{m s}^{-1}\). Its acceleration is \(-3\,\text{m s}^{-2}\) until the instant its velocity becomes zero. Immediately after this instant, its acceleration changes to \(+1\,\text{m s}^{-2}\) and remains constant. The position of the particle at \(t=5\,\text{s}\) is \hfill [4 marks]
{\renewcommand{\arraystretch}{1.8}
\begin{tabularx}{\linewidth}{X X X X}
(A) \quad \(-7.5\,\text{m}\) & (B) \quad \(6.0\,\text{m}\) & (C) \quad \(10.5\,\text{m}\) & (D) \quad \(17.5\,\text{m}\) \\[10pt]
\end{tabularx}}
\item A particle moves along a straight line with its initial position taken as the origin and the initial direction of motion taken as positive. It starts from rest and accelerates uniformly at $2\,\mathrm{m\,s^{-2}}$ for $5\,\mathrm{s}$. It then moves with a uniform deceleration of $5\,\mathrm{m\,s^{-2}}$ until it comes to rest. Subsequently, it reverses its direction and accelerates uniformly at $1\,\mathrm{m\,s^{-2}}$ in the negative direction. The reversed motion is considered up to the instant when its speed again becomes equal to the maximum speed attained before reversal. What is the displacement of the particle from its starting point when, during the reversed motion, its speed becomes one-fourth of the maximum speed attained in the entire motion? \hfill [4 marks]
{\renewcommand{\arraystretch}{1.8}
\begin{tabularx}{\linewidth}{X X}
(A) \quad $\dfrac{245}{8}\,\mathrm{m}$ & (B) \quad $\dfrac{255}{8}\,\mathrm{m}$ \\[10pt]
(C) \quad $\dfrac{265}{8}\,\mathrm{m}$ & (D) \quad $\dfrac{275}{8}\,\mathrm{m}$ \\[10pt]
\end{tabularx}}
\item A particle moves along a straight road, and its position relative to a fixed origin is given by $x(t)=t^3-6t^2+9t+2$ m for $0\leq t\leq 4$ s. The ratio of the total distance travelled by the particle to the magnitude of its net displacement during this interval is \hfill [4 marks]
{\renewcommand{\arraystretch}{1.8}
\begin{tabularx}{\linewidth}{X X X X}
(A) \quad $1:3$ & (B) \quad $2:1$ & (C) \quad $3:1$ & (D) \quad $4:1$ \\[10pt]
\end{tabularx}}
\item Two particles A and B move along the same straight track in the positive x-direction. At $t=0$, particle A is at $x=0\,\text{m}$ with velocity $8\,\text{m s}^{-1}$ and constant acceleration $2\,\text{m s}^{-2}$, while particle B is at $x=30\,\text{m}$ with velocity $20\,\text{m s}^{-1}$ and constant acceleration $-1\,\text{m s}^{-2}$. The position at which they meet for the first time after $t=0$ is: \hfill [4 marks]
{\renewcommand{\arraystretch}{1.8}
\begin{tabularx}{\linewidth}{X X X X}
(A) \quad $120\,\text{m}$ & (B) \quad $160\,\text{m}$ & (C) \quad $180\,\text{m}$ & (D) \quad $200\,\text{m}$ \\[10pt]
\end{tabularx}}
\end{enumerate}
\section*{Kinematics in One Dimension › Vertical motion under gravity}
\begin{enumerate}[resume]
\item An elevator is moving upward with a constant speed of \(5\,\text{m s}^{-1}\). At an instant when the elevator is \(30\,\text{m}\) above the ground, a passenger throws a ball vertically upward with a speed of \(15\,\text{m s}^{-1}\) relative to the elevator. Neglect air resistance and take \(g=10\,\text{m s}^{-2}\). What is the maximum height attained by the ball above the ground? \hfill [4 marks]
{\renewcommand{\arraystretch}{1.8}
\begin{tabularx}{\linewidth}{X X X X}
(A) \quad \(40\,\text{m}\) & (B) \quad \(45\,\text{m}\) & (C) \quad \(50\,\text{m}\) & (D) \quad \(55\,\text{m}\) \\[10pt]
\end{tabularx}}
\item A ball is projected vertically upward from the roof of a building. Its velocity--time graph is a straight line of slope $-g$ and crosses the time axis once; however, the portions corresponding to upward and downward motion are not marked. At two instants on this graph, the velocities are found to be $+18\,\text{m s}^{-1}$ and $-12\,\text{m s}^{-1}$, respectively. Taking $g=10\,\text{m s}^{-2}$, the time interval between these two instants is: \hfill [4 marks]
{\renewcommand{\arraystretch}{1.8}
\begin{tabularx}{\linewidth}{X X X X}
(A) \quad $1.8\,\text{s}$ & (B) \quad $2.4\,\text{s}$ & (C) \quad $3.0\,\text{s}$ & (D) \quad $4.2\,\text{s}$ \\[10pt]
\end{tabularx}}
\item An elevator is moving vertically upward with a constant acceleration of $2\,\mathrm{m\,s^{-2}}$. At a certain instant, a small ball is released from rest relative to the elevator at a height of $2.16\,\mathrm{m}$ above the elevator floor. Neglecting air resistance and taking $g=10\,\mathrm{m\,s^{-2}}$, the ball reaches the floor after \hfill [1 marks]
{\renewcommand{\arraystretch}{1.8}
\begin{tabularx}{\linewidth}{X X X X}
(A) \quad $0.4\,\mathrm{s}$ & (B) \quad $0.6\,\mathrm{s}$ & (C) \quad $0.8\,\mathrm{s}$ & (D) \quad $1.0\,\mathrm{s}$ \\[10pt]
\end{tabularx}}
\item A ball is projected vertically upward from the ground with an unknown initial speed. Its velocity–time graph is a straight line having slope $-10\,\text{m s}^{-2}$. At $t=1\,\text{s}$, its velocity is $+15\,\text{m s}^{-1}$. What is the total distance travelled by the ball during the interval from $t=1\,\text{s}$ to $t=4\,\text{s}$? \hfill [4 marks]
{\renewcommand{\arraystretch}{1.8}
\begin{tabularx}{\linewidth}{X X X X}
(A) \quad $15\,\text{m}$ & (B) \quad $22.5\,\text{m}$ & (C) \quad $30\,\text{m}$ & (D) \quad $45\,\text{m}$ \\[10pt]
\end{tabularx}}
\item A lift is moving vertically upward with a uniform acceleration of $2\,\mathrm{m\,s^{-2}}$. At a certain instant, a coin is released from rest relative to the lift from a point $2.4\,\mathrm{m}$ above the floor. Neglect air resistance and take $g=10\,\mathrm{m\,s^{-2}}$. The time taken by the coin to reach the floor, as measured by an observer inside the lift, is \hfill [4 marks]
{\renewcommand{\arraystretch}{1.8}
\begin{tabularx}{\linewidth}{X X X X}
(A) \quad $0.45\,\mathrm{s}$ & (B) \quad $0.63\,\mathrm{s}$ & (C) \quad $0.70\,\mathrm{s}$ & (D) \quad $0.89\,\mathrm{s}$ \\[10pt]
\end{tabularx}}
\item A ball is released from rest relative to an elevator that is moving upward with constant velocity. For the first $2\,\text{s}$ after release, the elevator continues with constant velocity. At the end of these $2\,\text{s}$, the elevator starts accelerating upward at $2\,\text{m s}^{-2}$. The ball then moves freely under gravity. Taking $g=10\,\text{m s}^{-2}$ and neglecting air resistance, which of the following correctly describes the motion of the ball relative to the elevator after the elevator starts accelerating? In particular, when will the ball again be at the same height as its release point in the elevator? \hfill [1 marks]
\begin{enumerate}[label=(\Alph*),leftmargin=1.6em,itemsep=1pt]\item After $\dfrac{10}{3}\,\text{s}$ \item After $\dfrac{5}{3}\,\text{s}$ \item After $2\,\text{s}$ \item It never returns to the release height for any positive time\end{enumerate}
\item An elevator is moving vertically upward with a constant speed of \(8\,\mathrm{m\,s^{-1}}\). At a certain instant, a small object is released from the elevator. Neglecting air resistance and taking \(g=10\,\mathrm{m\,s^{-2}}\), what is the separation between the object and the elevator after \(1.2\,\mathrm{s}\)? \hfill [4 marks]
\begin{enumerate}[label=(\Alph*),leftmargin=1.6em,itemsep=1pt]\item \(3.6\,\mathrm{m}\), with the object below the elevator \item \(7.2\,\mathrm{m}\), with the object below the elevator \item \(9.6\,\mathrm{m}\), with the object above the elevator \item \(14.4\,\mathrm{m}\), with the object below the elevator\end{enumerate}
\item An elevator is accelerating vertically upward with a constant acceleration of $2\,\text{m s}^{-2}$. At some instant, a ball is projected vertically upward with a speed of $4\,\text{m s}^{-1}$ relative to the elevator. Taking $g=10\,\text{m s}^{-2}$ and neglecting air resistance, after what time will the ball return to the point of projection as observed from the elevator? \hfill [4 marks]
{\renewcommand{\arraystretch}{1.8}
\begin{tabularx}{\linewidth}{X X X X}
(A) \quad $\frac{1}{3}\,\text{s}$ & (B) \quad $\frac{2}{3}\,\text{s}$ & (C) \quad $1\,\text{s}$ & (D) \quad $\frac{4}{3}\,\text{s}$ \\[10pt]
\end{tabularx}}
\item A lift is moving vertically upward with a constant acceleration of $2\,\text{m s}^{-2}$. A passenger inside the lift throws a small ball vertically upward with a speed of $12\,\text{m s}^{-1}$ relative to the lift. The passenger catches the ball when it returns to the same relative height. Neglect air resistance and take $g=10\,\text{m s}^{-2}$. The time interval between throwing and catching, as measured by the passenger, is \hfill [4 marks]
{\renewcommand{\arraystretch}{1.8}
\begin{tabularx}{\linewidth}{X X X X}
(A) \quad $1\,\text{s}$ & (B) \quad $2\,\text{s}$ & (C) \quad $2.4\,\text{s}$ & (D) \quad $3\,\text{s}$ \\[10pt]
\end{tabularx}}
\item A lift is moving vertically upward with a constant speed of \(5\,\text{m s}^{-1}\). When the lift is \(20\,\text{m}\) above the ground, a ball is projected vertically upward relative to the lift with speed \(15\,\text{m s}^{-1}\). Neglect air resistance and take \(g=10\,\text{m s}^{-2}\). What is the maximum height attained by the ball above the ground? \hfill [4 marks]
{\renewcommand{\arraystretch}{1.8}
\begin{tabularx}{\linewidth}{X X X X}
(A) \quad \(25\,\text{m}\) & (B) \quad \(35\,\text{m}\) & (C) \quad \(40\,\text{m}\) & (D) \quad \(45\,\text{m}\) \\[10pt]
\end{tabularx}}
\end{enumerate}
\end{document}